\documentclass{article}
\usepackage[T1]{fontenc}
\usepackage{lmodern}
\usepackage{graphicx}
\usepackage{amsmath,amssymb}
\usepackage[a4paper,margin=2cm]{geometry}
\usepackage[absolute,overlay]{textpos}
\setlength{\TPHorizModule}{1mm}
\setlength{\TPVertModule}{1mm}
\usepackage[indonesian]{babel}
\usepackage{hyperref}
\setlength{\emergencystretch}{2em}
% unit: Halaman 31

\begin{document}

% === Halaman 31 (llm) ===

Demikian seterusnya, sehingga plainteks dan cipherteks hasilnya adalah:
\begin{align*}
\text{Pesan (plainteks):} & \quad A23A9 \\
\text{Cipherteks (mode $ECB$):} & \quad 23124 \\
\text{Cipherteks (mode $CBC$):} & \quad 27FBF
\end{align*}

\end{document}
